% TOP_PROBLEM: {"id": 30006819, "title": "The Singer Conjecture on Vanishing of L2-Cohomology of Aspherical Manifolds", "category_id": 7, "category": {"id": 7, "name": "topology", "display_name": "Topology", "description": "Properties preserved under continuous deformations.", "slug": "topology", "order_index": 7, "created_at": "2026-07-31T15:26:25.671Z"}, "status": "open"}
% FIRST_POSTED: 2026-09-27
% ATTEMPT_STATUS: unresolved
% !TeX program = lualatex
\documentclass[11pt]{article}
\usepackage[margin=25mm]{geometry}
\usepackage{fontspec}
\setmainfont{Latin Modern Roman}
\usepackage{amsmath,amssymb,mathtools}
\usepackage{unicode-math}
\setmathfont{Latin Modern Math}
\usepackage{xurl,hyperref}
\hypersetup{colorlinks=true,urlcolor=blue}
\setcounter{secnumdepth}{0}
\setlength{\parindent}{0pt}
\setlength{\parskip}{5pt}
\setlength{\emergencystretch}{3em}
\begin{document}
\section*{\texorpdfstring{The Singer Conjecture on Vanishing of L2-Cohomology of Aspherical Manifolds}{The Singer Conjecture on Vanishing of L2-Cohomology of Aspherical Manifolds}}\label{30006819}
\subsection{Definitions and mathematical statement}
Let $M$ be a closed aspherical $n$-manifold with $G=\pi_1(M)$ and contractible universal cover $\widetilde M$. The $L^2$-Betti number $b_i^{(2)}(M)$ is the von Neumann dimension, over the group von Neumann algebra, of reduced $L^2$-homology of the lifted chain complex; on a smooth manifold it also measures the corresponding $L^2$ harmonic forms. The Singer assertion is
\[
 b_i^{(2)}(M)=0\qquad\text{whenever }i\ne n/2.
\]
For odd $n$, there is no integral middle degree, so every $L^2$-Betti number must vanish. These are invariants of the universal cover with its group action, not the ordinary Betti numbers of $M$.
\par\smallskip\textit{Formulation and concept references:} \href{https://link.springer.com/article/10.1007/s00039-024-00667-w}{[S1]}.

\subsection{Short English statement}
Must the L2-cohomology of a closed aspherical manifold be concentrated in its middle dimension, and vanish entirely in odd dimensions?
\subsection{Research attempt}
\paragraph{Scope, attribution, and source distinction.}
GPT-6 Astra Ultra prepared this unresolved attempt on 27 September 2026
using the formulation paper [S1], L\"uck's primary exposition [E1], and
explicit Hilbert-chain calculations. The target uses reduced $L^2$
homology and von Neumann dimension. In particular, the positive-characteristic
homology-growth phenomena discussed in [S1] are not automatically
counterexamples to this characteristic-zero $L^2$ assertion. The route
below proves the torus calculation directly, reconstructs the surface
case and product consequences from stated standard identities, and
tests why ordinary contractibility and Euler characteristic do not
supply the missing general vanishing.

\paragraph{Reduced homology is represented by a Hilbert-space kernel.}
For a bounded Hilbert chain complex with differential
$d_i:C_i\to C_{i-1}$, define
\[
 \overline H_i=\ker d_i/\overline{\operatorname{im}d_{i+1}},
 \qquad
 \Delta_i=d_i^*d_i+d_{i+1}d_{i+1}^*.
\]
The quotient identifies isometrically with
$\ker d_i\cap\ker d_{i+1}^*=\ker\Delta_i$.
Indeed the orthogonal complement of the closed boundary space is
$\ker d_{i+1}^*$, and
\[
 \langle\Delta_i v,v\rangle
       =\|d_iv\|^2+\|d_{i+1}^*v\|^2.
\]
This gives a concrete way to prove reduced vanishing without asserting
that the range of a differential is closed. The von Neumann dimension
of this equivariant harmonic space is the relevant Betti number for
the finite-cell models used below.

For example, let $S$ be the bilateral shift on $\ell^2(\mathbb Z)$.
The differential $S-I$ of the lifted circle complex is injective,
since its kernel consists of constant sequences and no nonzero constant
sequence is square summable. Its range is dense because the adjoint
kernel is also zero. Its range is not closed: the normalized indicator
of $\{1,\ldots,N\}$ has norm one while its image has norm
$\sqrt{2/N}$. If an injective operator had closed range, its inverse
on that range would be bounded, contradicting this sequence.
Thus reduced homology vanishes even though ordinary algebraic quotient
by the image need not vanish. Removing the closure from the catalogue
definition would change this simplest example.

\paragraph{An explicit Fourier calculation for every torus.}
For $n\ge1$, the universal cover of $\mathbb T^n$ has the standard cubical chain
complex with
$C_i=\ell^2(\mathbb Z^n)\otimes\bigwedge^i\mathbb C^n$.
Fourier transform identifies $\ell^2(\mathbb Z^n)$ with
$L^2(\mathbb T^n)$, where each deck shift becomes multiplication by
$z_j$. The differential is the Koszul contraction by the vector of
symbols $(z_1-1,\ldots,z_n-1)$, with the usual alternating signs.
The exterior contraction and wedge identities give
\[
                  \Delta_i(z)=
       \left(\sum_{j=1}^n|z_j-1|^2\right)I
\]
on every exterior degree. The multiplier is strictly positive except
at the single point $(1,\ldots,1)$, which has Haar measure zero.
An $L^2$ vector in its kernel vanishes almost everywhere. Therefore
\[
                   b_i^{(2)}(\mathbb T^n)=0
                   \quad\text{for all }i.
\]
This is compatible with the Singer assertion, including in even
dimension where the middle Betti number is permitted to be zero.

The inverse multiplier is unbounded near $(1,\ldots,1)$. Consequently
this calculation does not supply a bounded contracting homotopy of the
whole Hilbert complex, and does not justify an arbitrary interchange
between algebraic exactness and topological exactness. At the same
time, the ordinary Betti numbers of $\mathbb T^n$ are
$\binom ni$, often nonzero in many degrees. Ordinary cohomology of the
compact quotient is not what the conjecture asks to vanish.

\paragraph{The closed surface case from two structural identities.}
We use the established $L^2$ Euler formula and Poincar\'e duality from
[E1]: for a closed connected orientable finite-cell $n$-manifold,
\[
 \chi(M)=\sum_i(-1)^ib_i^{(2)}(M),
 \qquad b_i^{(2)}(M)=b_{n-i}^{(2)}(M).
\]
The zero-th Betti number is zero when the fundamental group is infinite.
One way to see the latter is that an invariant harmonic zero-chain on
a connected universal-cover graph is constant, and an infinite vertex
orbit admits no nonzero square-summable constant. These statements
refer to reduced homology, as in the preceding calculation.

For a closed orientable surface $\Sigma_g$ of genus $g\ge1$, its
fundamental group is infinite and its universal cover is contractible.
Thus $b_0^{(2)}=b_2^{(2)}=0$, and Euler characteristic gives
\[
                         b_1^{(2)}(\Sigma_g)=2g-2.
\]
This proves concentration in the middle degree in this family. For
nonorientable closed aspherical surfaces one can use the orientable
double cover and the established finite-cover scaling formula
$b_i^{(2)}(\widehat M)=[\pi_1M:\pi_1\widehat M]b_i^{(2)}(M)$,
also recorded in [E1]. No orientability assumption is silently added
to the original conjecture.

There is a useful ordinary-cover check. A degree-$d$ cover of
$\Sigma_g$ has genus $1+d(g-1)$, so its normalized first ordinary
Betti number is $2g-2+2/d$. The normalized zero-th and second Betti
numbers are $1/d$. Along degrees tending to infinity these match the
computed $L^2$ values. This arithmetic agreement in surfaces is not
a proof that every family of covers of every manifold has a corresponding
approximation property over every coefficient field.

\paragraph{Products preserve the predicted concentration.}
For finite Hilbert chain models of two spaces, the product differential
has the tensor-product signs. Its Laplacian on bidegree $(i,j)$ is
\[
                    \Delta_i^X\otimes I+I\otimes\Delta_j^Y.
\]
Both summands are nonnegative. Their sum has kernel precisely the
tensor product of their kernels: the spectral theorem reduces the
statement to $a+b=0$ with $a,b\ge0$, or equivalently one can use the
two quadratic forms and their zero spaces. Von Neumann dimensions
multiply for the product group. This proves the familiar $L^2$
K\"unneth formula in this setting,
\[
 b_k^{(2)}(X\times Y)=\sum_{i+j=k}b_i^{(2)}(X)b_j^{(2)}(Y).
\]
It follows that products of $r$ aspherical orientable surfaces have
only their degree-$r$ Betti number possibly nonzero, equal to
$\prod_{a=1}^r(2g_a-2)$. Their real dimension is $2r$, so this is
exactly the predicted middle degree. Multiplying by any positive-
dimensional torus makes all $L^2$-Betti numbers zero. In particular
these constructions verify many odd-dimensional cases as well.

More generally, if two closed aspherical factors satisfy Singer,
K\"unneth gives it for their product: two even-dimensional middle
degrees add to the middle degree, and an odd-dimensional factor has
all its Betti numbers zero. This closure observation does not express
an arbitrary aspherical manifold as such a product.

\paragraph{Why contractibility of the cover is insufficient.}
Asphericity kills ordinary reduced homology of the universal cover,
computed with finite chains or unrestricted ordinary cochains as
appropriate. It does not say that the necessary primitives are square
integrable. The genus-$g\ge2$ surface already provides a countertest:
its universal cover is contractible, but its first $L^2$-Betti number
is $2g-2>0$. That nonzero value is allowed by Singer because it occurs
in the middle. It proves that a proposed argument asserting all
$L^2$-cohomology vanishes merely from contractibility is wrong even
on a class for which the conjecture is known.

A contraction of a cell complex can involve chains moving arbitrarily
far or containing arbitrarily many translates. It need not extend as
a bounded operator on the square-summable completion. In the torus
calculation the unbounded inverse multiplier displays the same issue
analytically. A valid proof in a general manifold must control the
Hilbert complex itself, not just invoke an ordinary null-homotopy.

\paragraph{Euler sign is a consequence, not an equivalent reformulation.}
In even dimension $n=2m$, Singer implies
\[
                     (-1)^m\chi(M)=b_m^{(2)}(M)\ge0.
\]
In odd dimension it implies $\chi(M)=0$, but ordinary Poincar\'e
duality already yields that Euler equality. The reverse implications
are much weaker. As a purely numerical test, the symmetric list
$(b_0,b_1,b_2,b_3,b_4)=(0,1,4,1,0)$ has Euler sum two and respects
Poincar\'e symmetry, yet is not concentrated in degree two. In dimension
three the list $(0,1,1,0)$ has Euler sum zero but does not vanish.
These lists are not asserted to be realized by closed aspherical
manifolds. They show that nonnegativity, symmetry and the Euler sum
alone have not forced the individual off-middle terms to disappear.

The source [S1] further distinguishes rational $L^2$ phenomena from
mod-$p$ homology growth. Changing coefficients can introduce torsion
contributions invisible to rational dimensions. A mod-$p$ growth
counterexample must therefore not be relabeled as an $L^2$ counterexample
without a valid theorem identifying those invariants in that setting.

\paragraph{Verification, first gap, and outcome.}
Exact exterior-algebra checks verify the torus Laplacian symbol; the
surface and product formulas check dimensions, Euler signs and cover
normalizations. The structural duality and dimension facts are explicit
external inputs. The first missing general step is an analytic or
geometric mechanism forcing each off-middle harmonic Hilbert module
to vanish for arbitrary closed aspherical manifolds. No bounded
contraction, curvature hypothesis valid in every case, or universal
spectral estimate is obtained. The stated torus, surface and product
calculations are complete, while the general Singer conjecture remains
\textbf{UNRESOLVED}.

\subsection{Sources}
\begin{itemize}
\item[S1]Avramidi, Okun, and Schreve, Homology Growth, Hyperbolization, and Fibering, Geometric and Functional Analysis 34 (2024), Section 1.8.
\url{https://link.springer.com/article/10.1007/s00039-024-00667-w}.
\textit{Formulation source.}
\item[E1]Wolfgang L\"uck, \emph{$L^2$-Invariants from the Algebraic Point of View}, primary author exposition.
\url{https://him-lueck.uni-bonn.de/data/ltwoalg.pdf}.
\textit{Consulted 27 September 2026 for reduced Hilbert homology, von Neumann dimension, duality, Euler characteristics and the product formula.}
\end{itemize}
\end{document}
